
%\needspace{4\baselineskip}
\color{uniblue}{\subsection{Направленная защита от однофазных замыканий на землю (НЗОЗЗ)}}
\color{black}

\begin{enumerate}[label=\arabic{section}.\arabic{subsection}.\arabic{enumi}, labelsep=4pt, leftmargin=0pt, itemindent=57pt, itemsep=0pt, parsep=5pt] % если что можно крутить itemsep
%%%%%%%%%%%%%%%%%%%%%%%%%%%%%%%%%%%%%%%%%%%%%%%%%%%%%%%%%%%  ОБЩИЕ СВЕДЕНИЯ %%%%%%%%%%%%%%%%%%%%%%%%%%%%%%%%%%%%%%%%%%%%%%%%

\item Общие сведения

\begin{enumerate}[label=\arabic{section}.\arabic{subsection}.\arabic{enumi}.\arabic*, labelsep=4pt, leftmargin=0em, itemindent=65pt, parsep=0pt]

\item 
Направленная защита от однофазных замыканий на землю служит для селективного отключения однофазных замыканий на землю в сетях с изолированной нейтралью, а также в сетях с заземлением нейтрали через дугогасящий реактор или высокоомный резистор. В устройстве защита от однофазных замыканий на землю реализована с возможностью задания направленности.

\item 
В устройстве направленная защита от однофазных замыканий на землю представлена функциональным блоком «НЗОЗЗ». Основные характеристики функционального блока:
\begin{itemize}
    \item основная погрешность по углу при номинальных значениях тока и напряжения не более 5$^\circ$; 
    \item основная погрешность срабатывания по току, равному сумме высших гармоник, не более 20\%.
\end{itemize}

\item
Логическая часть реализует группу пусковых органов с независимыми уставками, предусматривает таймеры выдержек времени на срабатывание, логику пусков и блокировок, коммутацию внутренних пусковых сигналов для формирования выходного сигнала автоматического ускорения.

\item
В составе функционального блока «НЗОЗЗ» предусмотрены следующие функции:
\begin{itemize}
\item пусковой орган по 3I0 основной частоты;
\item пусковой орган по 3I0 высших гармоник;
\item пусковой орган по 3U0;
\item логика обнаружения переходных/перемежающихся замыканий;
\item орган направления мощности НП;
\item выходная логика.
\end{itemize}

\item
Функциональный блок «НЗОЗЗ» может быть полностью выведен из работы оперативно путем активации сигнала «ОВ НЗОЗЗ».
\item
В таблице \ref{nzozz:tbl1} приведены параметры, необходимые для настройки НЗОЗЗ.

\end{enumerate}

%%%%%%%%%%%%%%%%%%%%%%%%%%%%%%%%%%%%%%%%%%%%%%%%%%%%%%%%% Пусковой орган по 3I0 основной частоты (ПО 3I0) %%%%%%%%%%%%%%%%%%%%%%%%%%%%%%%%%%%%%%%%%%%%%%%%

\item Пусковой орган по току $3I_0$ основной частоты (ПО 3I0)

\begin{enumerate}[label=\arabic{section}.\arabic{subsection}.\arabic{enumi}.\arabic*, labelsep=4pt, leftmargin=0em, itemindent=65pt, parsep=0pt]


\item
В функциональном блоке «НЗОЗЗ» пусковой орган по току $3I_0$ основной частоты представлен функцией «ПО 3I0». Работа функции основана на сравнении значения первой гармоники тока нулевой последовательности с уставкой. Функция может использовать для работы либо непосредственно измеренный ток нулевой последовательности либо вычисленный (путем суммирования фазных токов), выбор осуществляется в ПО «ЮНИТ-Сервис» (вкладка «Конфигурация» $\rightarrow$  «Расчет 3I0» $\rightarrow$  «Измеренный» или «Вычисленный»).

\item
Ввод функции в работу осуществляется с помощью уставки SGF1 при введенной в работу функции выходной логики НЗОЗЗ.

\item
ПО предусматривает работу в трех режимах и регулируется уставкой SGF2:
\begin{itemize}
\item если уставка в положении «Ненаправленная», то работа «ПО 3I0» предусматривается без контроля направления даже при увеличении тока НП выше уставки;
\item если уставка в положении «Прямонаправленная», то работа «ПО 3I0» предусматривается только в прямом направлении;
\item если уставка в положении «Обратнонаправленная», то работа «ПО 3I0» предусматривается только в обратном направлении.
\end{itemize}

\item\sloppy
Работа функции при неисправности цепей напряжения определяется уставкой SGF3 «Реж\_БНН\_3U0» («Не предусмотрено», «Блокировка» или «Деблокировка»).

\item
В режиме, когда уставка находится в положение «Не предусмотрено», пусковые условия формируются без учета контроля цепей БНН в любом направлении. В режиме блокировки работа ступени запрещается до момента исчезновения неисправности БНН в любом направлении. В режиме деблокировки наличие неисправности в цепях напряжения, контролируемое по входному сигналу «БННш: Срабатывание», позволяет ПО функционировать без учета направления.

\item
Алгоритм работы функции выполнен в соответствии с рисунком \ref{nzozz:img1}. 

\vspace{1mm}
\begin{figure}[H]
\centering
\includegraphics[width=0.85\textwidth,height=0.85\textheight,keepaspectratio]{\fbpath/107. НЗОЗЗ/_latex/img/po3i0.pdf}
\captionof{figure}{Функциональная схема логики ПО 3I0}\label{nzozz:img1}
\end{figure}


\end{enumerate}

%%%%%%%%%%%%%%%%%%%%%%%%%%%%%%%%%%%%%%%%%%%%%%%%%%%%%%%%% Пусковой орган по 3I0 высших гармоник (ПО 3I0вг) %%%%%%%%%%%%%%%%%%%%%%%%%%%%%%%%%%%%%%%%%%%%%%%%

\item Пусковой орган по току $3I_0$ высших гармоник (ПО 3I0вг)

\begin{enumerate}[label=\arabic{section}.\arabic{subsection}.\arabic{enumi}.\arabic*, labelsep=4pt, leftmargin=0em, itemindent=65pt, parsep=0pt]

\item 
В функциональном блоке «НЗОЗЗ» пусковой орган по току $3I_0$ высших гармоник представлен функцией «ПО 3I0вг». Ввод функции в работу осуществляется с помощью уставки SGF1 при введенной в работу функции выходной логики НЗОЗЗ.

\item
В качестве воздействующей величины функция использует значение суммы высших гармонических составляющих (3, 5, 7, 9 и 11 гармоник) в токе нулевой последовательности. Полученное значение сравнивается с уставкой. Пуск функции осуществляется по истечению фиксированной выдержки времени таймера Т1.

% ИНФОРМАЦИЯ - 
\needspace{3\baselineskip}
\par
\noindent
\begin{tabularx}{\linewidth}{p{2em} X}
  \begin{minipage}[c][3\baselineskip][b]{2em}
    \centering\textcolor{uniblue}{\Huge\faCircleInfo}
  \end{minipage} &
  \hyphenpenalty=0 \textbf{При отсутствии в измерительном модуле канала для непосредственного измерения тока $3I_0$, работа пускового органа «ПО 3I0вг» не предусмотрена.}
\end{tabularx}
\vskip 2mm


\item
Алгоритм работы функции выполнен в соответствии с рисунком \ref{nzozz:img2}.

\vspace{1mm}
\begin{figure}[H]
\centering
\includegraphics[width=0.9\textwidth,height=0.9\textheight,keepaspectratio]{\fbpath/107. НЗОЗЗ/_latex/img/po3i0vg.pdf}
\captionof{figure}{Функциональная схема логики ПО 3I0вг}\label{nzozz:img2}
\end{figure}

\end{enumerate}

%%%%%%%%%%%%%%%%%%%%%%%%%%%%%%%%%%%%%%%%%%%%%%%%%%%%%%%%% Пусковой орган по 3U0 (ПО 3U0) %%%%%%%%%%%%%%%%%%%%%%%%%%%%%%%%%%%%%%%%%%%%%%%%

\item Пусковой орган по напряжению $3U_0$ (ПО 3U0)

\begin{enumerate}[label=\arabic{section}.\arabic{subsection}.\arabic{enumi}.\arabic*, labelsep=4pt, leftmargin=0em, itemindent=65pt, parsep=0pt]

\item 
В функциональном блоке «НЗОЗЗ» пусковой орган по напряжению $3U_0$ представлен функцией «ПО 3U0». Работа функции основана на сравнении первой гармоники напряжения нулевой последовательности с уставкой. Напряжение нулевой последовательности $3U_0$ получается путем непосредственного измерения или вычисления как сумма фазных напряжений.

\item
Ввод «ПО 3U0» в работу осуществляется с помощью уставки SGF1 и введенной в работу функции выходной логики НЗОЗЗ.

\item
Пусковые условия формируются при срабатывании «ПО 3U0» выше заданной уставки и отсутствие неисправности в цепях напряжения (БНН).

\item
Формирование разрешения от пускового органа осуществляется в зависимости от выбранной уставки программного переключателя, когда уставка SGF2 «Реж\_БНН\_3U0» находится в режиме деблокировки с активным сигналом от неисправности БНН или в режиме блокировки с отсутствием сигнала от БНН и срабатывании «ПО 3U0». В режиме, когда уставка находится в положении «Не предусмотрено», разрешение ПО формируется без учета контроля цепей БНН, пока сохранялись условия срабатывания функции.

\item
Алгоритм работы функции выполнен в соответствии с рисунком \ref{nzozz:img3}.

\vspace{1mm}
\begin{figure}[H]
\centering
\includegraphics[width=0.95\textwidth,height=0.95\textheight,keepaspectratio]{\fbpath/107. НЗОЗЗ/_latex/img/po3u0.pdf}
\captionof{figure}{Функциональная схема логики ПО 3U0}\label{nzozz:img3}
\end{figure}

\end{enumerate}
%%%%%%%%%%%%%%%%%%%%%%%%%%%%%%%%%%%%%%%%%%%%%%%%%%%%%%%%%%%%%%%%%%%%%%%%%%%%%%%%%%%%%%%%%%%%%%%%%%%%%%%%%%%%%%%%%%

%%%%%%%%%%%%%%%%%%%%%%%%%%%%%%%%%%%%%%%%%%%%%%%%%%%%%%%%% Логика обнаружения переходных/ перемежающихся замыканий (ЛОППЗ) %%%%%%%%%%%%%%%%%%%%%%%%%%%%%%%%%%%%%%%%%%%%%%%%

\item Логика обнаружения переходных/перемежающихся замыканий (ЛОППЗ)

\begin{enumerate}[label=\arabic{section}.\arabic{subsection}.\arabic{enumi}.\arabic*, labelsep=4pt, leftmargin=0em, itemindent=65pt, parsep=0pt]

\item
Основными воздействующими величинами являются мгновенные значения тока $3I_0$ и напряжения $3U_0$, полученные на частоте дискретизации 5 кГц.

\item В функциональном блоке «НЗОЗЗ» логика обнаружения переходных/перемежающихся замыканий представлена функцией «ЛОППЗ». Ввод функции в работу осуществляется с помощью уставки SGF1 при введенной в работу функции выходной логики НЗОЗЗ.

\item
В зависимости от выбранного положения уставки, возможны следующие режимы работы:
\begin{itemize}
    \item если уставка SGF1 выставлена в положение «По переходным процессам», то логика работы функции основана на контроле переходных процессов (см. пп. \ref{desc:transient});
    \item если уставка SGF1 выставлена в положение «По перемежающимся», то работа функции основана на контроле перемежающихся замыканий (см. пп. \ref{desc:fault}).
\end{itemize}

\item\label{desc:transient}
Защита от замыканий на землю по переходным процессам реализована путем измерения остаточных токов и напряжений и анализа переходных процессов тока $3I_0$ и напряжения $3U_0$ по заданным критериям. Измерение выполняется с высокой частотой дискретизации -- 5 кГц. Рабочими параметрами для функции являются мгновенные значения $3U_0$ и $3I_0$ (с частотой выборки 200 мкс, со знаком) и действующее значение $3U_0$.

\item
Данный режим рекомендуется к применению в качестве селективной защиты от замыканий на землю в сетях 6-35 кВ с компенсированной или изолированной нейтралью. Принцип действия основан на выявлении переходных процессов при замыканиях на землю. Если замыкание является устойчивым (не перемежающимся) по своей природе, то переходный процесс будет существовать только в начальный момент повреждения и это важно для измерения.

\item
Функция производит сравнение измеренных пиковых значений напряжения нулевой последовательности и тока нулевой последовательности с заданными пусковыми уровнями. При превышении заданного уровня производится определение переходных процессов в сигналах тока и напряжения нулевой последовательности. Имеется заранее заданный критерий для определения переходных процессов в сигналах $3U_0$ и $3I_0$  и их направления. Далее функция проверяет соответствие обнаруженных переходных процессов критериям направленности, установленных в параметрах режима направления.

\item
Выбор направления работы осуществляется уставкой SGF2 «Направление»:
\begin{itemize}
    \item если положение уставки выставлено в положение «Прямое», то условие для срабатывания может быть только, если направление переходных процессов как для $3I_0$, так и для $3U_0$ -- положительное или отрицательное;
    \item если положение уставки выставлено в положение «Обратное», то условие для срабатывания может быть только, если направления переходных процессов для $3U_0$ и $3I_0$ неодинаковы (одно положительное, другое отрицательное).
\end{itemize}

\item
Пусковые условия функции формируются при выполнении следующих условий:
\begin{itemize}
    \item число $n$ последовательных выборок $3U_0 > 3U_{0cp}$ равно уставке $N_{cp}$;
    \item число $n$ последовательных выборок $3I_0 > 3I_{0cp}$ равно уставке $N_{cp}$;
    \item орган определения направления передает разрешающий сигнал (например, когда направления обоих переходных процессов $3U_0$ и $3I_0$ одинаковы при уставке «Прямое», или различны при уставке «Обратное»).
\end{itemize}

\item
Счетчик $K$ необходим для подсчета количества импульсов («клевков»). Если количество импульсов превысило уставку, происходит пуск функции и срабатывание спустя время, задаваемое уставкой Т1.
\item
При восстановлении симметричного режима происходит сброс счетчика через регулируемую выдержку времени на возврат задаваемую уставкой Т3.

\item\label{desc:fault}
Режим работы по выявлению перемежающихся замыканий предназначен специально для определения проблем, вызванных ухудшением изоляции высоковольтных кабелей и особенно подходит для защиты кабельных сетей с компенсированной нейтралью.
\item
Его предназначение необходимо для определения замыканий вызванных:
\begin{itemize}
    \item ухудшением изоляции кабеля;
    \item механическим повреждением кабеля;
    \item старением изолирующего материала кабеля под действием попадающей на него влаги.
\end{itemize}
Все это может привести к снижению изоляционной стойкости и вызвать серию пробоев изоляции кабеля. Примерный вид импульса тока при замыканиях на землю показан на рисунке \ref{nzozz:transient}.

\item
В данном режиме используются мгновенные значения тока $3I_0$ и напряжения $3U_0$ с частотой выборки 5 кГц.

\item
В данном режиме производится сравнение измеренных пиковых значений напряжения нулевой последовательности и тока нулевой последовательности с заданными пусковыми уровнями. При превышении заданного уровня производится определение переходных процессов в сигналах тока и напряжения нулевой последовательности.

\item
Имеется заранее заданный критерий для определения переходных процессов в сигналах $3U_0$ и $3I_0$ и их направления. Далее проверяется соответствие обнаруженных переходных процессов критериям направленности, установленных в параметрах режима направления.

\vspace{1mm}
\begin{figure}[H]
\centering
\includegraphics[width=0.8\textwidth,height=0.8\textheight,keepaspectratio]{\fbpath/107. НЗОЗЗ/_latex/img/transient.jpg}
\captionof{figure}{Импульс тока при замыкании на землю}\label{nzozz:transient}
\end{figure}

\item
В качестве примера на рисунке \ref{nzozz:transient2} приведен график изменения тока и напряжения нулевой последовательности для поврежденного фидера и фидера без повреждения в сети с компенсированной нейтралью.

\vspace{1mm}
\begin{figure}[H]
\centering
\includegraphics[width=0.75\textwidth,height=0.75\textheight,keepaspectratio]{\fbpath/107. НЗОЗЗ/_latex/img/transient2.jpg}
\captionof{figure}{Ток нулевой последовательности поврежденного и неповрежденного фидера и напряжения нулевой последовательности при неустойчивых замыканиях на землю в сети с компенсированной нейтралью}\label{nzozz:transient2}
\end{figure}

\item
Когда сигнал «ЛОППЗ: Пуск» активен, запускается таймер T2. Если новый импульс происходит раньше окончания задержки на возврат ИО, таймер Т2 регистрирует один импульс.

\item
При достижении таймером T1 установленной выдержки времени на срабатывание, и по меньшей мере один импульс детектируется в течение задержки на возврат таймера Т2, активируется выходной сигнал «ЛОППЗ: Срабатывание». Сигнал будет активным, пока временной интервал между импульсами не станет больше задержки на возврат таймера Т2.

\item
Счетчик N необходим для подсчета количества выборок N. Если количество выборок превысило уставку происходит пуск функции и срабатывание спустя время, задаваемое уставкой Т1.

\item
При снижении тока возврата НП ниже уставки, по истечении 10 мс проиходит сброс счетчика числа выбросов, определяющих импульс.

\item
Алгоритм работы функции выполнен в соответствии с рисунком \ref{nzozz:img4}.

\vspace{1mm}
\begin{figure}[H]
\centering
\includegraphics[width=1\textwidth,height=1\textheight,keepaspectratio]{\fbpath/107. НЗОЗЗ/_latex/img/loppz.pdf}
\captionof{figure}{Функциональная схема логики ЛОППЗ}\label{nzozz:img4}
\end{figure}

\end{enumerate}
%%%%%%%%%%%%%%%%%%%%%%%%%%%%%%%%%%%%%%%%%%%%%%%%%%%%%%%%%%%%%%%%%%%%%%%%%%%%%%%%%%%%%%%%%%%%%%%%%%%%%%%%%%%%%%%%%%%%%%%%%%%%%%%%%%%%%%%%%%%%%%%%

%%%%%%%%%%%%%%%%%%%%%%%%%%%%%%%%%%%%%%%%%%%%%%%%%%%%%%%%% ОРГАН ОНМ НП %%%%%%%%%%%%%%%%%%%%%%%%%%%%%%%%%%%%%%%%%%%%%%%%

\item Орган направления мощности НП (ОНМ НП)

\begin{enumerate}[label=\arabic{section}.\arabic{subsection}.\arabic{enumi}.\arabic*, labelsep=4pt, leftmargin=0em, itemindent=65pt, parsep=0pt]

\item
В радиальных сетях, когда собственные емкостные токи отдельных присоединений велики и соизмеримы с полным током сети ненаправленная токовая защита нулевой последовательности не может обеспечить селективность действия. Для селективной работы в таких случаях необходимо применение органа направления мощности нулевой последовательности совместно с измерительным токовым органом защиты.

\item 
В функциональном блоке «НЗОЗЗ» орган направления мощности НП представлен функцией «ОНМ НП». Функция становится активной, если выходная логика НЗОЗЗ введена в работу с выставленным режимом работы («По 3I0» или «По 3I0 и 3U0») и настроен контроль направленности пускового органа $3I_0$ основной частоты.

\item
В качестве воздействующих величин функция использует значения первой гармоники тока нулевой последовательности и напряжения нулевой последовательности.

\item
Направленный элемент срабатывает при выполнении следующих условий:
\begin{itemize}
    \item напряжение НП превышает 1 вольт;
    \item ток НП превышает значение $0,04 \cdot I_{\textsl{ном}}$ НП;
    \item вектор полной мощности нулевой последовательности, определяемый по углу током НП и напряжением НП, находится в области срабатывания (прямое или обратное направление).
\end{itemize}

\item
Для определения области срабатывания задается уставка $\varphi_{\textsl{м.ч.}}$. Отсчет угла максимальной чувствительности производится от вектора напряжения НП и считается положительным, если направление отсчета совпадает с направлением вращения против часовой стрелки. Угол сектора отсчитывается от угла максимальной чувствительности в обе стороны и составляет 85°.

\item 
Сигнал «ОНМ НП: Прямое» появляется, если вектор полной мощности НП находится в диапазоне от $\varphi_{\textsl{м.ч.}}$ - 85° до  $\varphi_{\textsl{м.ч.}}$ + 85°. Сигнал «ОНМ НП: Обратное» появляется, если вектор полной мощности НП находится в диапазоне от ($\varphi_{\textsl{м.ч.}}$ - 180°) + 85° до ($\varphi_{\textsl{м.ч.}}$ - 180°) - 85°. Если вектор полной мощности НП находится в зоне «нечувствительности», то ни один из вышеуказанных сигналов не появляется.

\item 
Поясняющая диаграмма определения направления мощности приведена на рисунке \ref{nzozz:imgonm}.

\vspace{1mm}
\begin{figure}[H]
\centering
\includegraphics[width=0.6\textwidth,height=0.6\textheight,keepaspectratio]{\fbpath/107. НЗОЗЗ/_latex/img/char.pdf}
\captionof{figure}{Характеристика органа направления мощности нулевой последовательности}\label{nzozz:imgonm}
\end{figure}

\item
Алгоритм работы функции выполнен в соответствии с рисунком \ref{nzozz:imgonm1}.

\vspace{1mm}
\begin{figure}[H]
\centering
\includegraphics[width=0.9\textwidth,height=0.9\textheight,keepaspectratio]{\fbpath/107. НЗОЗЗ/_latex/img/onmnp.pdf}
\captionof{figure}{Функциональная схема логики ОНМ НП}\label{nzozz:imgonm1}
\end{figure}

\end{enumerate}

%%%%%%%%%%%%%%%%%%%%%%%%%%%%%%%%%%%%%%%%%%%%%%%%%%%%%%%%%%%%%%%%%%%%%%%%%%%%%%%%%%%%%%%%%%%%%%%%%%%%%%%%%%%%%%%%%%%%%%%%%%%%%%%%%%%%%%%%%%%

%%%%%%%%%%%%%%%%%%%%%%%%%%%%%%%%%%%%%%%%%%%%%%%%%%%%%%%%% Выходная логика НЗОЗЗ %%%%%%%%%%%%%%%%%%%%%%%%%%%%%%%%%%%%%%%%%%%%%%%%

\item Выходная логика НЗОЗЗ

\begin{enumerate}[label=\arabic{section}.\arabic{subsection}.\arabic{enumi}.\arabic*, labelsep=4pt, leftmargin=0em, itemindent=65pt, parsep=0pt]

\item
Выходные сигналы данной функции воздействуют на цепи отключения энергообъекта при выявлении аварийных режимов, выявленных функциями в составе НЗОЗЗ.

\item 
Ввод выходной логики в работу на этапе параметрирования устройства осуществляется программным переключателем SGF1 «Ввод\_функции». Информация о введенной выходной логики отображается выходным сигналом «ВыхЛогНЗОЗЗ: Ввод».

\item
Выходная логика НЗОЗЗ может быть выведена из работы оперативно путем активации сигнала «ОВ НЗОЗЗ», что характеризуется активным состоянием сигнала «ВыхЛогНЗОЗЗ: Оперативный вывод» на выходе функции.

\item
В зависимости от положения уставки SGF2, в функции возможны следующие режимы работы:
\begin{itemize}
    \item «По 3I0», то работа функции выполняется только по току $3I_0$ 50 Гц, расчет зависимой характеристики срабатывания осуществляется по току $3I_0$ 50 Гц;
    \item «По 3I0 и 3U0», то работа функции выполняется по току $3I_0$ 50 Гц и по напряжению $3U_0$, расчет зависимой характеристики срабатывания осуществляется по току $3I_0$ 50 Гц;
    \item «По 3I0вг», то работа функции выполняется только по току $3I_0$ высш. гарм., расчет зависимой характеристики срабатывания осуществляется по току $3I_0$ высш. гарм.;
    \item «По 3I0вг и 3U0», то работа функции выполняется по току $3I_0$ высш. гарм. и по напряжению $3U_0$, расчет зависимой характеристики срабатывания осуществляется по току $3I_0$ высш. гарм.
    \item «По 3U0», то работа функции выполняется только по пуску напряжения $3U_0$.
\end{itemize}

\item\sloppy
Контроль обнаружения перемежающихся замыканий вводится в работу уставкой SGF3 «Откл\_от\_ЛОППЗ» и действует на срабатывание функции без выдержки времени.

\item
При условии, что ключ «НЗОЗЗ на сигн» введен, по истечению независимой выдержки времени таймера Т2 действие функции переводится на сигнал.

\item
Выбор характеристик срабатывания для формирования сигнала  «ВыхЛогНЗОЗЗ: Срабатывание» и «ВыхЛогНЗОЗЗ: Сигнал» осуществляется программными переключателями «Характ\_ср» и «Характ\_сигн» соответственно. Графики реализованных в устройстве обратнозависимых времятоковых характеристик приведены в приложении \ref{app:chars}.

\item
Алгоритм работы выходной логики НЗОЗЗ выполнен в соответствии с рисунком \ref{nzozz:imgvl}. 

\vspace{1mm}
\begin{figure}[H]
\centering
\includegraphics[width=1\textwidth,height=1\textheight,keepaspectratio]{\fbpath/107. НЗОЗЗ/_latex/img/vl.pdf}
\captionof{figure}{Функциональная схема выходной логики НЗОЗЗ}\label{nzozz:imgvl}
\end{figure}

\end{enumerate}
%%%%%%%%%%%%%%%%%%%%%%%%%%%%%%%%%%%%%%%%%%%%%%%%%%%%%%%%% КОНЕЦ Выходная логика НЗОЗЗ %%%%%%%%%%%%%%%%%%%%%%%%%%%%%%%%%%%%%%%%%%%%%%%%

\end{enumerate}